\section*{Capítulo 7}

\solucion{7.1} $s_2=(1-0.9155)/(0.9615+0.9155)=0.0845/1.8770=4.50\,\%$.

\solucion{7.2} Valor para quien paga fijo $\approx(s'-s)\times\text{nocional}\times\sum DF=(0.0430-0.0480)\times50\times2.737=-0.68$ millones. La empresa se comprometió a pagar una tasa mayor que la que hoy conseguiría: el swap es un pasivo para ella.

\solucion{7.3} $s_{2,3}=\dfrac{DF(2)-DF(3)}{DF(3)}=\dfrac{0.911796-0.868608}{0.868608}=\dfrac{0.043188}{0.868608}=4.97\,\%$. Es la tasa a un año implícita entre los años 2 y 3.

\solucion{7.4} Por barril, PETROPERÚ recibe $89.38-80.00=9.38$. Para 1.16 millones de barriles: $10.88$ millones de dólares a su favor. El ingreso total por barril es lo que obtiene al vender el crudo (80.00) más la liquidación del swap (9.38): 89.38, el precio fijado.

\solucion{7.5} Porque las cuotas del préstamo pagan 5.55\,\%, más que la tasa que hoy exige el mercado en soles para esos plazos (alrededor de 4.0\,\%--4.5\,\%). Descontados a la curva, los pagos valen más que el saldo: S/ 3,838.4 millones frente a S/ 3,765.6 millones (US\$ 1,116.9 frente a 1,095.7 millones). Como el CCS debe valer cero al inicio, la pata en dólares tiene que valer también 1,116.9 millones sobre un nocional de 1,095.7 millones, lo que exige una tasa en dólares más alta.

\solucion{7.6} Porque toda la deuda ya está a tasa fija: no hay flujos variables que cubrir. Convertirla a variable añadiría riesgo de flujo de caja sin reducir el costo, que proviene sobre todo del spread de crédito (unos 569 pb) y no de la tasa base. El riesgo relevante es el de refinanciación, que se cubre mejor con un swap de inicio diferido o un \emph{rate lock} cuando se conozcan monto y fecha de la refinanciación del CESCE.
